\section{范式演进回顾：从任务专用监督到生成式智能体}

医疗人工智能在 2015--2026 年的演进可划分为三个相互衔接阶段。训练目标经历了由标签驱动进入语言驱动，再由语言驱动进入工具与反馈驱动的两次跃迁；每次跃迁都为模型拓展了新的能力边界。

\begin{figure}[htbp]
  \centering
  \fbox{\parbox{0.86\textwidth}{
    \centering
    图 2.1 占位图：三阶段训练目标演化\\[0.5em]
    阶段一为任务专用监督学习，阶段二为基础模型与多模态融合，阶段三为生成式 AI 与智能体。
  }}
  \caption{医疗 AI 三阶段训练目标演化示意。}
  \label{fig:paradigm-evolution}
\end{figure}

\noindent\textit{注：} 对比学习成为阶段二多模态对齐组件，BERT 风格双向编码器在阶段三中作为检索召回端或专科编码器继续发挥作用。

\subsection{阶段一：任务专用监督学习（约 2015--2020）}

这一阶段主要依托高质量标注数据，使用 CNN 与 RNN 进行端到端监督训练。Gulshan 等以 12.8 万张视网膜图像训练 InceptionV3，在糖尿病视网膜病变筛查上达到与眼科医师相当的灵敏度与特异性；Esteva 等以 12.9 万张皮肤镜图像训练 GoogLeNet 完成 2,032 种皮肤病变的细粒度分类；McKinney 等以英国与美国两国筛查数据训练乳腺癌检测模型，在前瞻验证中减少假阳性 5.7\%、假阴性 9.4\%[3]。截至目前，美国 FDA 已批准超过 500 项 AI/ML 医疗器械软件，其中绝大多数都属于这一阶段的产物。

尽管后续范式迅速迭代，这一阶段的许多方法仍沿用至今。\textbf{对比学习}自 CLIP、ConVIRT 起逐渐取代纯监督，把影像--报告对当作弱监督信号使用，为阶段二的多模态对齐组件奠定了基础；BioBERT、ClinicalBERT 则将 \textbf{BERT 风格的双向编码器} 适配到医学文本，成为 GatorTron、Med-BERT 等后续模型的技术基线[9]；以 \textbf{固定测试集上报告 AUC} 为核心的回顾性验证范式也在这一时期被确立，并在后续多年中保持为医疗 AI 论文默认的评测协议。

\subsection{阶段二：基础模型与多模态融合（2020--2023）}

这一时期的研究重心由单一任务的端到端训练，转向先通过自监督预训练构建通用底座，再迁移至具体任务。Transformer、AlphaFold、CLIP 与 GPT-3 共同构成了这一阶段的基础设施：Transformer 把序列建模统一成可扩展至跨模态的通用架构；AlphaFold 证明了大规模自监督在生物分子结构上的强迁移能力；CLIP 把图像与文本的对比学习推升为可扩展的视觉表示学习目标；GPT-3 则首次展示出不更新权重也能通过提示完成跨任务推理的零样本和少样本能力[4]。

在此基础上，Acosta 等于 2022 年提出多模态生物医学 AI（multimodal biomedical AI）概念框架，把精准组学健康、数字化临床试验、远程监护、大流行病监测、数字孪生与虚拟健康助手六类应用纳入统一视角[5]；Moor 等于 2023 年提出通用医疗 AI（generalist medical AI, GMAI）框架，将通用医疗模型的三大核心能力概括为动态任务指定、灵活多模态输入输出与医学领域知识表示[6]。这两套框架共同确立了阶段二的目标函数，不再局限于在单一任务上达到医生水平，而是追求以同一组权重在更广泛任务谱系上保持可用性。

若干代表性体系也在这一阶段集中出现。语言侧，GatorTron 以 820 亿临床笔记词、60 亿 PubMed 词和 25 亿 Wiki 词从头训练至 89B 参数，在临床推理任务上比 BioBERT 提升约 9.5--9.6\%[9]；BioGPT 在 1,500 万 PubMed 摘要上从头训练 GPT-2 Medium，PubMedQA 准确率达到 78.2\%[10]。多模态侧，LLaVA-Med 通过 600K 图文对齐加 60K 自指令对话的两阶段课程，仅使用 8 张 A100、在不足 15 小时内完成训练，在生物医学视觉问答任务上达到与远大于自身的模型相当的性能[11]；Med-Flamingo 以 OpenFlamingo + LLaMA-7B 为骨架，构造从 4,721 本医学教科书抽取的 MTB 数据集（80 万图、5.84 亿 tokens），首次在医学多模态领域验证了少样本上下文学习的可行性[12]。

\subsection{阶段三：生成式 AI 与智能体（2023--2026）}

2022 年末 ChatGPT 引发广泛关注后，医疗 AI 的方法学焦点迅速由基础模型本身转向生成式接口、工具调用与推理增强的组合。Med-PaLM 系列是这一阶段的代表性起点。Singhal 等在 540B Flan-PaLM 上仅通过指令提示微调，更新约 65 个软提示参数，就把科学共识一致率从 61.9\% 提升到 92.6\%，接近临床医生的 92.9\%[8]；Med-PaLM 2 在 PaLM 2 基础上引入集成精炼（ensemble refinement）与检索链（chain of retrieval）两项推理增强机制，使 MedQA 准确率达到 86.5\%，并在 9 个临床轴中的 8 个上获得医生评审者偏好[13]。

Med-Gemini 进一步拓展了基于不确定性的自主搜索能力。当模型输出 token 的 Shannon 熵超过阈值时，系统自动触发网络检索，把外部知识接地纳入推理回路，同时利用 Gemini 1.5 的 1M+ token 上下文窗口，使模型能够直接解析完整电子病历或手术视频，无需依赖 RAG 切片[7,16]。AMIE 则把自博弈与模拟学习作为创新点：在 PaLM 2 基础上，内循环通过 in-context 批评反馈优化对话，外循环通过迭代微调更新权重，模拟覆盖 5,230 种疾病的多轮问诊场景，最终在 159 例 OSCE 风格双盲交叉研究中以 30/32 个专家维度与 25/26 个患者维度显著优于 20 名全科医生[17]。

该阶段的技术路径呈相互促进的发展态势。一方面，推理增强开始进入医疗领域，与传统的指令微调形成互补；另一方面，工具调用与智能体协议让模型从纯权重模型转向带工具的智能体。O'Sullivan 等于 2026 年发表的心脏专科 AMIE，以 Gemini 2.0 Flash 为底座，配套领域工具与多步推理，仅基于 9 例真实病例进行开发，即完成了 107 例 HCM 与心肌病的随机交叉部署，显著错误率较单独专家组下降近一半[19]。Teo 等在 2025 年的综述中将这一阶段的方法学谱系概括为：合成数据系统、扩散模型、规则型 AI、LLM、多模态基础模型、推理与智能体，以及模型蒸馏[7]。

\subsection{技术里程碑时间表}

表 \ref{tab:milestone} 给出 2017 年至 2026 年医疗 AI 关键里程碑的时间表。纵向阅读对应三阶段演进，横向阅读对应训练范式的多样化。

\begin{table}[htbp]
\centering
\caption{医疗 AI 关键里程碑（2017--2026）}
\label{tab:milestone}
\scriptsize
\begin{tabularx}{\textwidth}{L{1.6cm} L{4.5cm} L{2.3cm} Y}
\toprule
时间 & 里程碑 & 所属阶段 & 主要贡献 / 意义 \\
\midrule
2017 & Transformer[A] & 阶段二基础设施 & 跨模态统一架构基础 \\
2018--2019 & BERT / BioBERT / ClinicalBERT[A] & 阶段一$\rightarrow$二 & 生物医学双向编码器，迁移学习起点 \\
2020 & GPT-3 / CLIP / AlphaFold[4,A] & 阶段二基础设施 & LLM、视觉对比学习与蛋白结构三大基石 \\
2021 & Med-BERT[30] / CEHR-BERT[31] & 阶段一$\rightarrow$二 & 结构化 EHR 预训练，奠定专科基础模型路线 \\
2022 & GatorTron[9] / BioGPT[10] / Multimodal Biomedical AI[5] & 阶段二 & 域内规模定律首次验证，多模态综述框架成形 \\
2022 & ChatGPT[A] & 阶段二$\rightarrow$三 & 引爆生成式医疗 AI 范式 \\
2023 & GMAI[6] / Med-PaLM[8] / LLaVA-Med[11] / Med-Flamingo[12] / RETFound[32] & 阶段三起点 & 通用医疗 AI 框架、医学 LLM 与多模态 VLM 并行突破；首个眼科自监督基础模型 \\
2023 & HuatuoGPT[33] / BianQue[34] / DISC-MedLLM[35] / SAM-Med2D[36] & 阶段三 & 中文医疗 LLM 路线启航；分割基础模型扩展至 10 模态 \\
2024 & Med-PaLM 2[13] / Med-Gemini[7,16] / BiomedGPT[15] / MAIRA-2[24] / CONCH[37] / CT-FM (Pai)[38] / MedSAM[39] / EquityMedQA[40] & 阶段三 & 多模态、长上下文、不确定性引导并进；专科基础模型扩展至病理与 CT；公平性工具箱出现 \\
2024 & Foresight 2[41] / AMIE[17] & 阶段三 & EHR 时序专科 LLM；自博弈对话诊断 \\
2025 & MetaMedQA[2] / MedHELM[18] / FUTURE-AI[42] & 阶段三评估转向 & 元认知评估、整体性评估与国际可信 AI 共识 \\
2026 & 心脏专科 AMIE[19] / PreA RCT[20] / 心理治疗认知层[21] / Bean RCT[26] & 阶段三临床证据 & 真实临床随机对照试验证据集中涌现 \\
\bottomrule
\end{tabularx}
\end{table}

\noindent\textit{A：} 时间线背景节点引自 Acosta 等[5]、Moor 等[6]、Thirunavukarasu 等[22]、Teo 等[7]的综述以及通用 AI 文献基础。

\subsection{章节小结}

医疗人工智能经历了多代方法学的演进和叠加。阶段一的对比学习经 CLIP / ConVIRT 演化，成为阶段二多模态对齐的主流组件；阶段二的 BERT 风格双向编码器在阶段三中以检索召回端或专科编码器的形式继续发挥作用，Med-BERT、CEHR-BERT、RETFound、CONCH 等专科基础模型正是这一脉络在 2021--2024 年的延续；阶段一确立的回顾性验证范式则至今仍作为绝大多数医疗 AI 论文默认的评测协议在使用。换言之，医疗 AI 的范式更替并不是简单替代，而是前代技术在后续框架中的重组和再激活。
